\chapter{Resultados esperados e limitações}
\label{sec:conclusao}

Espera-se que a pesquisa entregue quatro resultados, aqui ordenados pelo que o
desenho tem mais condição de medir. O primeiro é a medida da aplicação aérea
na Chapada do Apodi nos anos anteriores à lei, que hoje não existe em fonte
pública. O segundo é a evidência sobre a substituição de tecnologia e as
intoxicações de trabalhadores. O terceiro é o ajuste produtivo: área,
produção e emprego, e a avaliação parcial de custo-benefício, em cenários. O
quarto é o efeito sobre os desfechos perinatais. O peso ao nascer continua o
desfecho primário, mas a média municipal funciona como limite superior: o
intervalo pode excluir efeitos grandes, de deriva ampla, e não distingue a
ausência de efeito do efeito localizado do cenário de referência
(Apêndice~\ref{ap:plano}). Da avaliação do desenho já sai uma recomendação que
não depende das estimativas: a política carece de monitoramento, e os
registros de aplicação e de receituário permitiriam instituí-lo.

O projeto tem limitações que convém declarar. A exposição da gestante não é
observada, e todos os efeitos são de intenção de tratar no nível do município.
Há dois municípios tratados no núcleo, três com Russas, e os intervalos serão condicionais a hipóteses sobre a dependência e a
representatividade dos choques, que os dados de controle não verificam por
inteiro. O primeiro estágio depende de pedidos de acesso à
informação, que podem ser negados ou devolver registros incompletos. A linha de
base do método de aplicação é de 2006. Um choque estadual simultâneo à lei só é
tratado pelas especificações secundárias, e o segundo contraste compara
sistemas produtivos distintos. O desenho não tem poder para detectar, na
média municipal, o efeito do cenário de referência. A análise do segundo
evento deve ficar descritiva. E a análise não é inteiramente cega: a
série estadual de produção e os números de dois municípios tratados são
públicos e foram vistos, e uma versão anterior do projeto, com outro desenho e
só com municípios cearenses, estimou efeitos sobre o peso ao nascer depois de
2019.
